%! Linear Functions: Slope, Forms \& Graphing — Unit 2, Lesson 2.1 — Group Activity (Answer Key)
\documentclass[10pt]{article}
\usepackage{algebra2-article}
\usepackage{algebra2-key}   % pulls in -boxes; do NOT also load -boxes
\newcommand{\TallMath}[1]{$\displaystyle #1\rule[-1.4em]{0pt}{3.2em}$}
\newcommand{\lgrid}[4]{%
  \draw[step=1, linegray!40, very thin] (#1,#3) grid (#2,#4);
  \draw[->, semithick, charcoal] (#1,0)--(#2,0) node[below right, font=\scriptsize]{$x$};
  \draw[->, semithick, charcoal] (0,#3)--(0,#4) node[left, font=\scriptsize]{$y$};}

\begin{document}

\pageheader{Unit 2, Lesson 2.1}{Group Activity --- Answer Key}
\namepartnerperiod

\vspace{0.08in}

\begin{headlinebox}{forestbg}
\small\textbf{Building Lines.} Yesterday you read a line's features off a graph; today you
\emph{build} the equation. Every line comes from a \textbf{slope} and \textbf{one point}. Work with
your partner and keep your forms tidy: integers in standard form, with $A \ge 0$.
\end{headlinebox}

\vspace{0.06in}

\begin{tcolorbox}[colback=white, colframe=black!40, title=\textbf{Tier R --- Remediate},
    fonttitle=\bfseries, arc=1.5mm, left=3mm, right=3mm, top=2mm, bottom=2mm, breakable]
Each equation is in slope-intercept form $y=mx+b$. Read off its slope and $y$-intercept, then match
it to one of the three graphs by writing \textbf{I}, \textbf{II}, or \textbf{III}.
\begin{center}
\begin{tikzpicture}[scale=0.4]
  \begin{scope}
    \lgrid{-2}{4}{-2}{4}
    \draw[very thick, forest] (-0.5,3.5)--(3.5,-0.5);
    \fill[charcoal] (0,3) circle (3pt); \fill[charcoal] (3,0) circle (3pt);
    \node[charcoal, font=\scriptsize] at (1,-3){Graph I};
  \end{scope}
  \begin{scope}[xshift=8cm]
    \lgrid{-2}{4}{-3}{4}
    \draw[very thick, forest] (-0.5,-2)--(2.5,4);
    \fill[charcoal] (0,-1) circle (3pt); \fill[charcoal] (1,1) circle (3pt);
    \node[charcoal, font=\scriptsize] at (1,-4){Graph II};
  \end{scope}
  \begin{scope}[xshift=16cm]
    \lgrid{-2}{4}{-2}{4}
    \draw[very thick, forest] (-2,0)--(4,3);
    \fill[charcoal] (0,1) circle (3pt); \fill[charcoal] (2,2) circle (3pt);
    \node[charcoal, font=\scriptsize] at (1,-3){Graph III};
  \end{scope}
\end{tikzpicture}
\end{center}
\renewcommand{\arraystretch}{1.6}
\begin{tabularx}{\linewidth}{@{}>{\bfseries}p{0.28\linewidth} X X X@{}}
 & \textbf{Slope $m$} & \textbf{$y$-intercept $(0,b)$} & \textbf{Graph} \\
\hline
$y = 2x - 1$        & \ans{$2$} & \ans{$(0,-1)$} & \ans{II} \\
$y = -x + 3$        & \ans{$-1$} & \ans{$(0,3)$} & \ans{I} \\
$y = \tfrac{1}{2}x + 1$ & \ans{$\tfrac{1}{2}$} & \ans{$(0,1)$} & \ans{III} \\
\end{tabularx}
\end{tcolorbox}

\begin{tcolorbox}[colback=white, colframe=black!40, title=\textbf{Tier A --- Approaching Proficiency},
    fonttitle=\bfseries, arc=1.5mm, left=3mm, right=3mm, top=2mm, bottom=2mm, breakable]
Write each line's equation. Show the point-slope step, then give slope-intercept \emph{and} standard
form.
\begin{enumerate}[leftmargin=1.7em, itemsep=10pt]
  \item \textbf{Slope and a point.} Slope $-2$, through $(3,1)$.\\
        Point-slope: \ans{$y - 1 = -2(x - 3)$} \\[10pt]
        Slope-intercept: \ans{$y = -2x + 7$} \qquad Standard: \ans{$2x + y = 7$}
  \item \textbf{Two points.} Through $(1,1)$ and $(3,7)$.\\
        Slope $m = \ans{$\dfrac{7-1}{3-1} = 3$}$ \\[8pt]
        Point-slope: \ans{$y - 1 = 3(x - 1)$} \\[10pt]
        Slope-intercept: \ans{$y = 3x - 2$} \qquad Standard: \ans{$3x - y = 2$}
  \item \textbf{Justify.} Are $y = 3x - 2$ and $3x - y = 2$ the \emph{same} line? Explain how you can
        tell without graphing.
        \ansline{Yes. Rearrange the standard form: $3x - y = 2 \Rightarrow -y = 2 - 3x \Rightarrow y = 3x - 2$. It becomes the slope-intercept equation exactly, so they name the same line.}
\end{enumerate}
\end{tcolorbox}

\begin{tcolorbox}[colback=white, colframe=black!40, title=\textbf{Tier E --- Extension},
    fonttitle=\bfseries, arc=1.5mm, left=3mm, right=3mm, top=2mm, bottom=2mm, breakable]
\begin{enumerate}[leftmargin=1.7em, itemsep=10pt]
  \item \textbf{Transform the parent.} Describe how to obtain the graph of $y = -3x + 2$ from the
        parent $y = x$, in terms of a \emph{stretch/reflection} and a \emph{shift}.
        \ansline{Reflect $y=x$ over the $x$-axis and stretch it vertically by a factor of $3$ (slope $1 \to -3$), then shift the graph up $2$ (the $y$-intercept moves from $(0,0)$ to $(0,2)$).}
  \item \textbf{Parallel \& perpendicular.} A line is given by $y = 2x - 4$.
        \begin{enumerate}[label=(\alph*), leftmargin=1.8em, itemsep=6pt]
          \item Write the line \emph{parallel} to it through $(0,1)$: \ans{$y = 2x + 1$}
          \item Write the line \emph{perpendicular} to it through $(0,1)$: \ans{$y = -\tfrac{1}{2}x + 1$}
        \end{enumerate}
  \item \textbf{In context.} Your streaming service costs $C(m) = 2m + 5$ dollars after $m$ months.
        A rival charges \$3 per month with no signup fee, $D(m) = 3m$.
        \begin{enumerate}[label=(\alph*), leftmargin=1.8em, itemsep=6pt]
          \item What do the slope and $y$-intercept of each model mean in words?
                \ansline{$C$: slope $2$ = \$2 per month; intercept $5$ = \$5 signup fee. $D$: slope $3$ = \$3 per month; intercept $0$ = no signup fee.}
          \item After how many months do the two services cost the \emph{same}, and what is that
                cost? Show your work.
                \ansline{Set $2m+5 = 3m \Rightarrow 5 = m$. At $m=5$ months, $D(5)=3(5)=\$15$ (and $C(5)=2(5)+5=\$15$). They match at $5$ months, \$15.}
        \end{enumerate}
\end{enumerate}
\end{tcolorbox}

\begin{teachernote}
Tier~A item~3 is the conceptual crux: ``same line, different form'' is verified by \emph{algebra},
not appearance. In Tier~E, watch the perpendicular slope --- students often negate \emph{or}
reciprocate but not both; $m=2 \Rightarrow m_\perp = -\tfrac12$. Item~3's break-even ties slope
(rate) and intercept (fixed cost) back to the streaming Hook; accept a clear table or an equation
for full credit.
\end{teachernote}

\end{document}
